\section{Results}

\subsection{Dataset composition}
\begin{figure}[tb]
  \centering
  \includegraphics[width=\textwidth]{fig1_dataset.pdf}
  \caption{\textbf{Target selection, the cliff definition, and two confounds it
    carries.} (a)~Six ChEMBL targets screened for activity cliffs, split into
    constitutional and stereochemical pairs. Only stereoisomer-free constitutional
    cliffs are a fair test of a 2D descriptor: no stereo-blind fingerprint can
    resolve a stereoisomer pair by construction, while any 3D method resolves it
    trivially, so a benchmark dominated by them favours structure before a single
    model is fitted. Acetylcholinesterase is the only target screened where
    constitutional cliffs dominate, 100 against 10; BACE1 inverts that at 5
    against 101. (b)~All 97 cliff pairs in the final set are constitutional and
    differ by 1--10 heavy atoms, yet binary ECFP4 reports Tanimoto exactly 1.000
    for 49 of them. Count-based ECFP4 separates the same pairs at a median of
    0.878. (c)~The cliff definition on the selected set; the similarity axis is
    the binary fingerprint, which is the criterion that selected the pairs.
    (d)~Cliff partners differ by 4.4 heavy atoms on average against 1.8 for
    controls.}
  \label{fig:dataset}
\end{figure}

\begin{table}[tb]
  \centering
  \caption{\textbf{Dataset composition.} The test split has not been scored by any
    model in this work.}
  \label{tab:dataset}
  \begin{tabular}{lrrrrr}
\toprule
Split & Ligands & of which cliff & Pairs & Cliff pairs & Control pairs \\
\midrule
Development & 217 & 81 & 145 & 77 & 68 \\
Test (sealed) & 54 & 18 & 38 & 20 & 18 \\
\midrule
Total & 271 & 99 & 183 & 97 & 86 \\
\bottomrule
\end{tabular}

\end{table}

\TODO{Prose. Composition from Table~\ref{tab:dataset}; the two confounds from
  Figure~\ref{fig:dataset}.}
\NOTE{The saturation in panel (b) is not only an observation about the pairs: the
  selection criterion is that same binary fingerprint, verified identical to the
  audit's t\_bit column. So the measure that chose the pairs is the one that
  cannot tell molecules eight heavy atoms apart from identical. Say that plainly
  -- it affects how the cliff set itself should be read, and it is a property of a
  default representation rather than of this target.}
\NOTE{Size confound: stage 10 found 0 of 19 descriptors survive a size adjustment,
  and four that looked significant are size artefacts.}

% ---------------------------------------------------------------------------
\subsection{The ladder}
\label{sec:ladder}

\begin{figure}[tb]
  \centering
  \includegraphics[width=\textwidth]{fig2\_ladder.pdf}
  \caption{\textbf{No structure-based rung resolvably outperforms 2D
    fingerprints.} Each point is the best score over a rung's descriptor arms and
    over the algorithm families, on 217 development ligands and 145 development
    pairs with frozen five-fold cross-validation. Rung~1 carries both feature sets
    benchmarked for it, because the winner depends on the framing. Rung~2 appears
    twice, as a single rigid receptor and as the median over eight receptor
    conformations. The dashed line marks the rung-1 level.
    \TODO{add interval description once fig\_ladder.py is run with --n-boot 300}}
  \label{fig:ladder}
\end{figure}

\begin{table}[tb]
  \centering
  \caption{\textbf{Best score per rung on each framing}, over the rung's descriptor
    arms and the algorithm families, on 217 development ligands and 145 development
    pairs. The winning family is named in parentheses. The full grid is
    the supplementary tables at the end of this section.}
  \label{tab:benchmark}
  \begin{tabular}{lccc}
\toprule
& \framingA{} & \framingB{} & \framingC{} \\
Rung & $R^2$ & $R^2$ & ROC AUC \\
\midrule
Rung 1 2D & +0.386 {\scriptsize (PLS)} & +0.200 {\scriptsize (Gradient Boosting)} & 0.805 {\scriptsize (Random Forest)} \\
Rung 2 pose & +0.046 {\scriptsize (PLS)} & +0.167 {\scriptsize (SVM (RBF))} & 0.684 {\scriptsize (Gradient Boosting)} \\
Rung 2 ensemble & +0.215 {\scriptsize (Random Forest)} & +0.274 {\scriptsize (SVM (RBF))} & 0.705 {\scriptsize (Random Forest)} \\
Rung 3 MM/GBSA & +0.118 {\scriptsize (Random Forest)} & -0.056 {\scriptsize (Random Forest)} & 0.753 {\scriptsize (Random Forest)} \\
Rung 4 MD & +0.068 {\scriptsize (Random Forest)} & +0.124 {\scriptsize (SVM (RBF))} & 0.793 {\scriptsize (Random Forest)} \\
\bottomrule
\end{tabular}

\end{table}

\TODO{Prose. The numbers, from results/20\_ladder.csv and Table~\ref{tab:benchmark}:}
\NOTE{\framingA{}: rung 1 $+0.386$, rung 2 $+0.046$, rung 2e $+0.215$, rung 3
  $+0.118$, rung 4 $+0.056$. \framingB{}: $+0.200$, $+0.167$, $+0.274$, $-0.056$,
  $+0.124$. \framingC{}: $0.805$, $0.684$, $0.705$, $0.753$, $0.793$.
  Rung 3 on framings A and C is new -- it had only ever been scored on B.
  All five rungs cover all 217 ligands and all 145 pairs, so no rung is scored on
  an easier subset.}

\TODO{State plainly that on \framingB{} the ensemble rung 2 is numerically above
  rung 1 ($+0.274$ vs $+0.200$) and that stage 15 showed the difference is not
  resolvable. Do not write that 2D wins everything.}

\subsection{Rung 4 against a matched 2D baseline}
\TODO{Gaps and intervals, 300-replicate bootstrap: \framingA{} $-0.330$,
  CI $[-0.607,-0.203]$ -- excludes zero, rung 4 is worse. \framingB{} $-0.077$,
  CI $[-0.428,+0.131]$. \framingC{} \TODO{pending}.}
\NOTE{Framing A's interval is the only one in the study that excludes zero. It
  does not say rung 4 is unproven; it says rung 4 is worse, resolved by the design.}

% ---------------------------------------------------------------------------
\subsection{What the study could have detected}

\begin{figure}[tb]
  \centering
  \includegraphics[width=\textwidth]{fig3\_power.pdf}
  \caption{\textbf{The design separates large effects only, and what it leaves
    unexplained is signal rather than assay noise.}
    (a)~Power to call a candidate descriptor set better than rung~1, with the
    baseline re-simulated inside every replicate; 80\% power requires
    $R^2 \approx 0.45$ against a rung-1 baseline of 0.209.
    (b)~The 95\% interval on the gap is 0.32 wide on average and nearly flat
    across the range.
    (c)~The measurement-noise ceiling, from 336 of 5{,}609 compounds carrying a
    repeat assay, is 0.616 even on the pessimistic p90 spread -- far above every
    rung.}
  \label{fig:power}
\end{figure}

\TODO{Prose. The portable sentence is: this study separates two $R^2$ values about
  0.30 apart and not less. Derived before any rung-4 run finished, and it held.}

% ---------------------------------------------------------------------------
\subsection{Descriptor reliability}
\begin{figure}[tb]
  \centering
  \includegraphics[width=\textwidth]{fig4_reliability.pdf}
  \caption{\textbf{Of three descriptor design classes, only persistence survives a
    replicate test.} (a)~Ratio of the spread across ligands to the spread across
    replicates of the same ligand, for each descriptor at each analysis window;
    the threshold of 2 was fixed before the pilot ran. (b)~Verdict composition by
    window. Not one of the 20 event-count descriptor--window combinations is
    usable at any length, against 57 of 100 for persistence.}
  \label{fig:reliability}
\end{figure}

\TODO{Prose. The counts: 81 of 160 combinations seed-dominated; by design class,
  persistence 57/100 usable, fluctuation 6/35, event \textbf{0/20}. At the 20~ns
  window 11 survive and 8 of them are contact occupancies.}
\NOTE{The zero is the quotable number. How long a contact persists reproduces
  across seeds; how fast something fluctuates, at this length, largely does not.}

\TODO{Show that the filter earns its place: \framingB{} $+0.124$ filtered against
  $-0.054$ unfiltered; \framingC{} $0.793$ against $0.721$.}

% ---------------------------------------------------------------------------
\subsection{Trajectory length}

\begin{figure}[tb]
  \centering
  \includegraphics[width=\textwidth]{fig5\_duration.pdf}
  \caption{\textbf{Trajectory length is not the limiting factor between 5 and
    20~ns.} Descriptors recomputed on the first 5, 10, 20 and 30~ns of analysed
    trajectory. The per-window replicate filter selects 15, 14 and 12 descriptors
    respectively, so it confounds trajectory length with descriptor choice; the
    fixed series holds the 20~ns selection at every window and isolates length.
    Only 45 development ligands reach 30~ns, drawn as their own series and never
    joined to the full-sample line.}
  \label{fig:duration}
\end{figure}

\TODO{Prose. With a fixed descriptor set, \framingA{} is flat
  ($+0.039$, $+0.027$, $+0.056$) and \framingB{} peaks at 10 ns
  ($+0.116$, $+0.198$, $+0.124$). The spread is inside the study's resolution, so
  the conclusion is no resolvable dependence on length -- which answers the cost
  objection with data.}
\NOTE{Do NOT quote the $+0.198$ as a result. It is the best of six cells computed
  together and it is the kind of number this paper exists to warn about.}

% ---------------------------------------------------------------------------
\subsection{A partial subset inflates an expensive descriptor set's apparent gain}
\label{sec:subset}

\begin{table}[tb]
  \centering
  \caption{\textbf{The gain reported at 59 of 145 pairs was an artefact worth
    $+0.277$.} Prefixes follow the real completion order, recovered from each
    run's file timestamp; random subsets of the same ligand count are the control.
    Gains are against rung~1's published full-set figure and against rung~1
    recomputed on the same pairs.}
  \label{tab:subset}
  \input{../results/21\_subset\_inflation.tex}
\end{table}

\TODO{Prose. Two findings: it is not a small-sample effect (random subsets of the
  same size put rung 4 at $+0.093$ where arrival order puts it at $+0.477$), and
  recomputing the baseline is necessary but not sufficient -- it removes 0.094 of
  the 0.277 and leaves $+0.183$ standing.}
\NOTE{Mechanism: state all three co-factors, not the convenient one. The early
  ligands are smaller (median 31 vs 35 heavy atoms, $p=0.0005$, decaying to
  $p=0.54$ by 196 ligands), they hold every multi-replicate run, and they average
  28.8 ns of analysed trajectory against 20.0. 217 ligands cannot separate them.}

% ---------------------------------------------------------------------------
\subsection*{Supplementary tables}
\TODO{Move to the SI document when one exists.}

\begin{table}[tb]
  \centering
  \caption{Every rung, descriptor arm and regression family on
    \framingA{} and \framingB{}.}
  \begin{tabular}{lllrr}
\toprule
Rung & Descriptor arm & Algorithm & \framingA{} $R^2$ & \framingB{} $R^2$ \\
\midrule
Rung 1 2D & binary ECFP4 & Ridge (linear) & +0.175 & +0.106 \\
 &  & PLS & +0.386 & +0.100 \\
 &  & k-NN (k=5) & +0.104 & +0.099 \\
 &  & SVM (RBF) & +0.383 & +0.124 \\
 &  & Random Forest & +0.329 & +0.184 \\
 &  & Gradient Boosting & +0.375 & +0.200 \\
\addlinespace
Rung 1 2D & count ECFP4 + RDKit2D & Ridge (linear) & +0.077 & -513.621 \\
 &  & PLS & +0.357 & -687.740 \\
 &  & k-NN (k=5) & +0.081 & -0.637 \\
 &  & SVM (RBF) & +0.380 & -0.519 \\
 &  & Random Forest & +0.359 & -0.328 \\
 &  & Gradient Boosting & +0.336 & -0.007 \\
\addlinespace
Rung 2 pose & 4EY7 + bridging & Ridge (linear) & +0.045 & -0.186 \\
 &  & PLS & +0.046 & -0.120 \\
 &  & k-NN (k=5) & -0.074 & -0.075 \\
 &  & SVM (RBF) & -0.213 & +0.167 \\
 &  & Random Forest & +0.031 & +0.063 \\
 &  & Gradient Boosting & -0.058 & +0.024 \\
\addlinespace
Rung 2 ensemble & ensemble median & Ridge (linear) & +0.183 & -0.029 \\
 &  & PLS & +0.131 & -0.018 \\
 &  & k-NN (k=5) & +0.155 & +0.134 \\
 &  & SVM (RBF) & +0.201 & +0.274 \\
 &  & Random Forest & +0.215 & +0.043 \\
 &  & Gradient Boosting & +0.182 & +0.190 \\
\addlinespace
Rung 3 MM/GBSA & MM/GBSA & Ridge (linear) & +0.034 & -0.183 \\
 &  & PLS & +0.026 & -0.259 \\
 &  & k-NN (k=5) & -0.045 & -0.239 \\
 &  & SVM (RBF) & -0.045 & -0.228 \\
 &  & Random Forest & +0.118 & -0.056 \\
 &  & Gradient Boosting & +0.041 & -0.110 \\
\addlinespace
Rung 4 MD & replicate-filtered (12) & Ridge (linear) & +0.013 & +0.084 \\
 &  & PLS & -0.028 & +0.072 \\
 &  & k-NN (k=5) & -0.153 & +0.012 \\
 &  & SVM (RBF) & -0.067 & +0.124 \\
 &  & Random Forest & +0.056 & -0.084 \\
 &  & Gradient Boosting & -0.028 & -0.104 \\
\addlinespace
Rung 4 MD & all descriptors (33) & Ridge (linear) & -0.186 & -0.148 \\
 &  & PLS & -0.279 & -0.054 \\
 &  & k-NN (k=5) & +0.014 & -0.060 \\
 &  & SVM (RBF) & -0.307 & -0.108 \\
 &  & Random Forest & +0.068 & -0.073 \\
 &  & Gradient Boosting & -0.110 & -0.166 \\
\addlinespace
\bottomrule
\end{tabular}

\end{table}

\begin{table}[tb]
  \centering
  \caption{Every rung, descriptor arm and classification family
    on \framingC{}. The families differ from Table~S1 because \framingC{} is a
    classification; Ridge and PLS have no counterpart here and are absent rather
    than failed.}
  \begin{tabular}{lllr}
\toprule
Rung & Descriptor arm & Algorithm & \framingC{} ROC AUC \\
\midrule
Rung 1 2D & binary ECFP4 & Logistic regression & 0.597 \\
 &  & k-NN (k=5) & 0.590 \\
 &  & SVM (RBF) & 0.657 \\
 &  & Random Forest & 0.648 \\
 &  & Gradient Boosting & 0.520 \\
\addlinespace
Rung 1 2D & count ECFP4 + RDKit2D & Logistic regression & 0.609 \\
 &  & k-NN (k=5) & 0.683 \\
 &  & SVM (RBF) & 0.606 \\
 &  & Random Forest & 0.805 \\
 &  & Gradient Boosting & 0.575 \\
\addlinespace
Rung 2 pose & 4EY7 + bridging & Logistic regression & 0.574 \\
 &  & k-NN (k=5) & 0.614 \\
 &  & SVM (RBF) & 0.555 \\
 &  & Random Forest & 0.602 \\
 &  & Gradient Boosting & 0.684 \\
\addlinespace
Rung 2 ensemble & ensemble median & Logistic regression & 0.639 \\
 &  & k-NN (k=5) & 0.681 \\
 &  & SVM (RBF) & 0.635 \\
 &  & Random Forest & 0.705 \\
 &  & Gradient Boosting & 0.673 \\
\addlinespace
Rung 3 MM/GBSA & MM/GBSA & Logistic regression & 0.689 \\
 &  & k-NN (k=5) & 0.698 \\
 &  & SVM (RBF) & 0.630 \\
 &  & Random Forest & 0.753 \\
 &  & Gradient Boosting & 0.655 \\
\addlinespace
Rung 4 MD & replicate-filtered (12) & Logistic regression & 0.749 \\
 &  & k-NN (k=5) & 0.737 \\
 &  & SVM (RBF) & 0.741 \\
 &  & Random Forest & 0.793 \\
 &  & Gradient Boosting & 0.785 \\
\addlinespace
Rung 4 MD & all descriptors (33) & Logistic regression & 0.627 \\
 &  & k-NN (k=5) & 0.587 \\
 &  & SVM (RBF) & 0.535 \\
 &  & Random Forest & 0.721 \\
 &  & Gradient Boosting & 0.700 \\
\addlinespace
\bottomrule
\end{tabular}

\end{table}
